\section{Introduction}
\label{introduction}

RNA-binding proteins (RBPs) regulate RNA splicing, stability, localization, and
translation. Determining where these proteins contact RNA is therefore essential
for studying post-transcriptional gene regulation. Enhanced crosslinking and
immunoprecipitation (eCLIP) provides transcriptome-wide measurements of such
contacts~\parencite{vannostrand2016}. However, the assay produces aligned
sequencing reads rather than a definitive list of binding sites. Peak-calling
software must distinguish genuine crosslink signal from experimental background,
and the resulting calls determine the quality of downstream analyses.

PureCLIP infers crosslink sites with a hidden Markov model and can incorporate a
size-matched input control~\parencite{krakau2017}. Its output depends on parameters
that control signal smoothing, site merging, and filtering. A parameter set that
works well for one experiment may not transfer to another because RBPs differ in
binding specificity and eCLIP libraries differ in depth, background, and antibody
performance. Manual selection is consequently difficult to standardize and
becomes impractical when each genome-wide evaluation requires substantial
compute time.

This work addresses that problem with an automated optimization loop. At each
iteration, an optimizer proposes a parameter set, a shared workflow runs
PureCLIP and post-processing, and the resulting sites are scored for replicate
reproducibility, motif support, and recovery of ENCODE reference regions. We
compare two proposal strategies: an LLM agent that receives decomposed feedback
and RBP-specific biological context, and the Tree-structured Parzen Estimator
(TPE) implemented in Optuna. Both strategies use the same evaluation code,
search bounds, and composite objective. A second LLM condition withholds the
biological context and thereby tests whether the prior changes search performance.

The study asks three questions. First, does automated search improve on the
default PureCLIP configuration? Second, does the LLM attain higher-quality call
sets than TPE under the available experimental protocol? Third, does RBP-specific
context benefit the LLM, and does that benefit vary with the sequence specificity
of the RBP? The software registry contains 26 ENCODE datasets covering 15 RBPs in
K562 and HepG2 cells. The analysis reported here uses a completed paired subset:
seven K562 and five HepG2 dataset--cell-line pairs representing eight RBPs. Six
K562 datasets also have an LLM-without-prior arm. The registry permits the same
pipeline to be run on additional proteins without changing the evaluation code,
but those datasets were not part of the available completed cohort. Conclusions
therefore concern this subset rather than the full registry.
